% 16 — Distributive law  A ∩ (B ∪ C) = (A ∩ B) ∪ (A ∩ C)
\begin{tikzpicture}[font=\itshape]
  \newcommand{\threeV}[2]{\begin{scope}[shift={(#1,#2)}]
      \coordinate (cA) at (0,0.65);\coordinate (cB) at (-0.7,-0.55);\coordinate (cC) at (0.7,-0.55);
      \def\r{1.05}
      \draw[gray,line width=.8pt,rounded corners=3pt] (-2.1,-2.1) rectangle (2.1,2.05);
      \node[gray,font=\scriptsize\upshape,below left] at (2.1,2.05) {$U$};
      \begin{scope}\clip (cA) circle(\r);
        \fill[orange,opacity=0.5] (cB) circle(\r) (cC) circle(\r);\end{scope}
      \draw[navy,line width=.9pt] (cA) circle(\r);
      \draw[Fred,line width=.9pt] (cB) circle(\r);
      \draw[Tgreen,line width=.9pt] (cC) circle(\r);
      \node[navy,font=\small\itshape] at ($(cA)+(0,1.25)$){$A$};
      \node[Fred,font=\small\itshape] at ($(cB)+(-1.05,-0.45)$){$B$};
      \node[Tgreen,font=\small\itshape] at ($(cC)+(1.05,-0.45)$){$C$};\end{scope}}
  \node[formula,font=\itshape\normalsize] at (-3.4,2.55) {$A \cap (B \cup C)$};
  \threeV{-3.4}{0}
  \node[Tgreen,font=\large\upshape] at (0,0.7) {\cmark};
  \node[navy,font=\LARGE\upshape] at (0,-0.2) {$=$};
  \node[formula,font=\itshape\normalsize] at (3.4,2.55) {$(A \cap B) \cup (A \cap C)$};
  \threeV{3.4}{0}
\end{tikzpicture}
